\chapter{Происхождение магнитного сверхтонкого члена}
\label{ch:v8-baryon-em-magnetic-hyperfine-origin-gate}

Пространственный гейт сохранил условный электростатический знак, но не
включал магнитные моменты. Здесь необходимо прежде всего разделить слабый
ароматный дублет и физический спин: матрицы Паули предыдущего
изоспинового расчёта действовали на аромат, а не на новый спиновый носитель.

\section{Минимальное расширение носителя}

Для одного кварка требуется пространство
\begin{equation}
 \mathcal H_{q}=\mathbb C^2_{\rm aroma}\otimes\mathbb C^2_{\rm spin},
 \qquad
 \mathcal H_{3q}=\mathcal H_q^{\otimes3}.
 \label{eq:v8-baryon-em-spin-flavor-carrier}
\end{equation}
Пусть $\boldsymbol\Sigma_r$ --- физические матрицы Паули на спиновом
множителе, а $Q_r$ --- заряд на ароматном множителе. Минимальный
контактный оператор имеет вид
\begin{equation}
 H_{\rm mag}=\frac{\zeta}{T}
 \sum_{1\le r<s\le3}h_{rs}Q_rQ_s\,
 \boldsymbol\Sigma_r\!\cdot\!\boldsymbol\Sigma_s.
 \label{eq:v8-baryon-em-magnetic-contact-operator}
\end{equation}
Коэффициент $\zeta$ содержит знак, магнитные моменты и размерную
нормировку, а $h_{rs}$ --- контактный пространственный матричный элемент.
Ни один из них не является следствием оператора слабого изоспина.

\section{Полного спина недостаточно}

Рассмотрим два нормированных состояния с $S=1/2$, $S_z=1/2$:
\begin{equation}
 \chi_0=\frac{|\uparrow\downarrow\uparrow\rangle
 -|\downarrow\uparrow\uparrow\rangle}{\sqrt2},
 \qquad
 \chi_1=\sqrt{\frac23}|\uparrow\uparrow\downarrow\rangle
 -\frac{|\uparrow\downarrow\uparrow\rangle
 +|\downarrow\uparrow\uparrow\rangle}{\sqrt6}.
 \label{eq:v8-baryon-em-two-spin-half-couplings}
\end{equation}
Их попарные корреляции равны соответственно
\begin{equation}
 \bigl(\langle\Sigma_1\!\cdot\!\Sigma_2\rangle,
 \langle\Sigma_1\!\cdot\!\Sigma_3\rangle,
 \langle\Sigma_2\!\cdot\!\Sigma_3\rangle\bigr)
 =(-3,0,0),\quad(1,-2,-2).
 \label{eq:v8-baryon-em-spin-half-pair-correlations}
\end{equation}
При $h_{rs}=1$ для фиксированных ароматных рисунков $uud$ и $udd$
оператор
$O=\sum_{r<s}Q_rQ_s\boldsymbol\Sigma_r\cdot\boldsymbol\Sigma_s$
даёт
\begin{equation}
 \begin{array}{c|cc|c}
  &O_p&O_n&O_n-O_p\\ \hline
  \chi_0&-4/3&2/3&2\\
  \chi_1& 4/3&0&-4/3
 \end{array}.
 \label{eq:v8-baryon-em-spin-coupling-sign-countermodels}
\end{equation}
Знак меняется при одном и том же полном спине. Эти состояния не объявляются
готовыми фермионными волновыми функциями: они являются точным
свидетельством того, что нужен дополнительный селектор совместной
перестановочной симметрии спина, аромата и пространства.

\section{Условное симметричное спин-ароматное состояние}

Обозначим через $P_{\rm sf}$ усреднение по одновременным перестановкам
трёх спин-ароматных множителей:
\begin{equation}
 P_{\rm sf}=\frac1{6}\sum_{\pi\in S_3}U_\pi^{\rm aroma}
 \otimes U_\pi^{\rm spin}.
 \label{eq:v8-baryon-em-spin-flavor-symmetrizer}
\end{equation}
В секторе $P_{\rm sf}=1$, $I=S=1/2$, $I_z=\pm1/2$, $S_z=1/2$
пространство для протона и нейтрона одномерно. На его нормированных
векторах точная зарядово-спиновая форма равна
\begin{equation}
 O_p=\frac43,
 \qquad
 O_n=1,
 \qquad
 O_n-O_p=-\frac13.
 \label{eq:v8-baryon-em-symmetric-spin-flavor-magnetic-values}
\end{equation}

\begin{proposition}[Условная магнитная разность]
Если принять полностью симметричное спин-ароматное состояние, общий
контакт $h_{rs}=h$ и оператор \eqref{eq:v8-baryon-em-magnetic-contact-operator},
то
\begin{equation}
 E_n^{\rm mag}-E_p^{\rm mag}=-\frac{\zeta h}{3T}.
 \label{eq:v8-baryon-em-conditional-magnetic-splitting}
\end{equation}
Алгебраический множитель $-1/3$ фиксирован, но физический знак остаётся
знаком произведения $-\zeta h$.
\end{proposition}

Полная симметрия спин-ароматной части обычно связывается с
антисимметричным цветом и симметричным основным пространственным состоянием.
В текущем проекте последний переход не выведен: эпсилон-проектор задаёт
цвет, но не создаёт физический спиновый множитель и не доказывает
пространственную $S$-волну.

\section{Контактный масштаб}

Для координатного растяжения предыдущей главы контактная плотность
масштабируется как
\begin{equation}
 \left\langle\delta^{(d)}(r_i-r_j)\right\rangle_{\psi_s}
 =s^d\left\langle\delta^{(d)}(r_i-r_j)\right\rangle_{\psi}.
 \label{eq:v8-baryon-em-contact-density-dilation}
\end{equation}
Следовательно, даже условно выбранная спин-ароматная волновая функция не
фиксирует магнитную величину без радиального родителя, масс кварков и
нормировки магнитных моментов.

\begin{nogocriterion}[Граница сверхтонкого члена]
Слабый изоспин $\mathbf I^2$ нельзя переименовывать в физический спин и
использовать как сверхтонкий гамильтониан. Для такого вывода нужны
независимый спиновый носитель, совместный перестановочный селектор,
пространственная контактная плотность и коэффициент магнитного момента.
\end{nogocriterion}

\begin{protocol}
\begin{enumerate}
 \item вычислительное поле: \textbf{точное, без чисел с плавающей точкой};
 \item ароматный и физический спиновые носители: \textbf{разделены};
 \item два состояния с $S=1/2$: \textbf{построены точно};
 \item знак из одного полного спина: \textbf{не следует};
 \item симметричный сектор $I=S=1/2$: \textbf{одномерен};
 \item условный множитель $O_n-O_p=-1/3$: \textbf{точен};
 \item знак и величина $\zeta h$: \textbf{не выведены};
 \item контактная масштабная орбита: \textbf{$s^d$};
 \item физическое магнитное расщепление: \textbf{не получено};
 \item следующий гейт: \textbf{совместный перестановочный селектор
 спина, аромата и пространства}.
\end{enumerate}
\end{protocol}

Машинный сертификат сохранён в
\path{s2t_v8_baryon_em_magnetic_hyperfine_origin_gate_results.json}.